\section{Motivation}

\subsection{Threat Model}

\subsubsection{The ontology identification problem}
We can't yet reliably look inside a model and draw conclusions about its
goals or world model. Interpretability tools have recovered a hidden
objective in an auditing game with a planted behavior, but a semantic
search of the training data surfaced the same evidence without inspecting
the model at all~\cite{marks2025}. This is a major limitation on almost
all proposed alignment agendas, because behavioral testing alone cannot
settle what a model is optimizing for: a model can behave differently
when it infers that it is being trained or evaluated, and when that
context comes from training documents rather than the prompt, the gap
can persist without any visible chain of thought~\cite{greenblatt2024}.
Although not every model exhibits this gap~\cite{sheshadri2025}, the
broader point that aligned behavior during testing need not predict
behavior after deployment is standard~\cite[Ch.~8]{bostrom2014}.

\paragraph{Inner misalignment}
We don't currently know how to specify a robust goal for a model.
Because goal space is large, we should expect
goal misgeneralization as the default outcome.


\subsubsection{Limits of natural abstractions}
Networks trained on the same task can converge on the same algorithm
across architectures while differing, even across random seeds, in which
features implement it~\cite{chughtai2023}. Similarly, in language models,
SAE features from a Transformer and a Mamba trained on the same corpus
match nearly as well as those from two Transformer seeds~\cite{wang2024}.
Both results support a weak form of the universality hypothesis---the
claim within the natural abstraction hypothesis that diverse cognitive
systems learn similar high-level abstractions~\cite{lang2023}---which,
if true, would provide a convenient shortcut for auditing models.

Natural abstractions could be considered the set of concepts we could use to communicate with aliens. 
They should consistently emerge through unsupervised learning (in the limit given some constraints).
If any neural network is considered a fuzzy generalization of a finite state machine, 
natural abstractions are the latent states.

Unfortunately, the abstractions a model learns seem highly sensitive to choice of training data, and even to the random seed with the data fixed~\cite{chughtai2023}, and are thus not necessarily universal in the limit.
\footnote{The same applies to those that seem most natural to humans.}
Assuming that model capabilities will continue to primarily be driven by access to more data and more modalities, we cannot expect the abstractions learned by a weak model to remain robust in a more powerful model.

Even if we could fully characterize a model's state space in terms of natural abstractions,
that still doesn't mean we can reconstruct the latent FSM, 
as the transition rules can be arbitrarily complex.


\subsection{Theory of change}
Conditional on the underlying model of ontology being true, the
techniques proposed here would be a substantial interpretability
advance. The agenda also has a case for counterfactual impact: its
theoretical basis consists of nonobvious applications of obscure
techniques from niche fields of bioinformatics, and similar research
is unlikely to be done soon otherwise.

\subsubsection{What a solution would look like}
Mechinterp plausibly offers the best path to a robust alignment
solution despite criticisms.
The case rests on
\begin{enumerate}
    \item \textbf{decomposability} Mechinterp allows alignment to be
    broken into manageable subproblems.
    \item \textbf{parallelizability} Decomposable subproblems can be
    worked on independently.
    \item \textbf{robustness} Mechinterp offers direct causal
    explanations of model behavior.
\end{enumerate}
The most straightforward mechinterp alignment solutions are
\begin{enumerate}
    \item eliciting latent knowledge (ELK)
    Directly inferring a model's ontology.
    \label{it:elk}
    \item mechanistic anomaly detection (MAD)
    Detecting when a model has ``unusual'' thoughts.
    \label{it:mad}
\end{enumerate}
Both of these can be considered
ontology identification subproblems. MAD requires mapping between the model's activations and
the model's ontology. ELK requires mapping between the model's ontology
and human ontology. These are hard but not intractable problems.
Recently, developmental interpretability
(devinterp) techniques have identified distinct ``developmental
stages'' during training, each marked by a shift in internal structure
and behavior such as learning bigram statistics or forming induction
heads~\cite{hoogland2024developmental}. This suggests, at least in small
transformers, that training consists largely of the network
grokking~\cite{nanda2023progress} discrete concepts. While these results
are promising for MAD, locating these emergent capabilities in the
model's activations remains an open problem.

Interpretability glosses broadly as ``let humans understand AI''. The
alternative---getting AI to understand humans better---scales badly:
making the model smarter leads to subtler and subtler issues with larger
and larger effects. Between the two sits scalable oversight, in which
models help humans supervise other models; humans paired with an
unreliable dialog assistant already outperform both the assistant and
unaided humans~\cite{bowman2022}, and interpretability provides another
signal such oversight can use. Moreover, what we're doing is not just
letting humans understand AI. Because an LLM reflects humanity from its
training data, high quality conceptual interpretability provides
machine-comprehensible mathematical representations of every concept
that humanity cares about. We work on interp for the usual
corrigibility reasons, and also because outer alignment needs a stable
target to attach values to: a goal defined in terms of a model's internal
concepts must continue to point at those concepts as the model's
representations are updated~\cite[Ch.~12]{bostrom2014}. This agenda
attacks both problems, though our preliminary results show how hard
concept stability is even at small scale: past the first layer, learned
heads do not survive a change of random seed
(Section~\ref{sec:prelim-seed}).

\subsubsection{Macro-scale strategy}
Broadly, the goals are to
\begin{enumerate}
  \item formalize the notion of a model's latent concept space
  \item distill a model's latent concept space
  \item enforce decomposability of latent concept space
  \item enforce robustness of existing concepts as the dimensionality of
  the latent concept space increases
\end{enumerate}
A prerequisite for all of these is the ability to express a
metaontology
which can identify latent concepts in activation space.


